\largerpage
\sscp{*ə/*a/*-ay/*-aw > **a /σ\und{σ}{\#} (final syllables)}{11-01-02}
The second change shared by STB languages is merger of *ə/*a/*-ay/*-aw
> **a in final syllables.\footnote{
		It is worth noting that both
		\ili{Fordata} and \ili{Kei} are described synchronically
		as having centralised allophones of /a/ in final syllables.
		Thus, for instance, \ili{Fordata} has /a/ → [ə] in final open syllables
		\citep[192]{Marshall2000}.}
The change *-ay > \it{a} is found
in all other languages of central Maluku (\srf{07-04-01}),
and *-aw > \it{a} occurs in nearly all languages of central Maluku,
except \ili{Alune} (\srf{10-02}).
*ə > \it{a} in final syllables is unique among languages
of central Maluku, but it is common in \tsc{Timor-Babar}.
In particular, the \tsc{Southwest Tanimbar} node neighbouring
STB has *ə > \it{a} (\srf{04-07}).
Likewise, *ə > \it{a} in final syllables also occurs in some
neighbouring SHWNG languages \citep[79f]{Kamholz2014}.

\begin{table}\tcp{\tsc{Seram-Tanimbar-Bomberai} *ə > **a /σσ{\#}}{11-04}
\begin{adjustbox}{width=\textwidth}
	\begin{threeparttable}\st{2pt}
	\begin{tabular}{llllllll}\lsptoprule
local gl.	&	black	&	tooth	&	hear	&	roof	&	rope	&	inside	&	six\\
PMP	&	*ma-qit\tbr{ə}m	&	*[n/l]ip\tbr{ə}n	&	*d\tbr{ə}ŋ\tbr{ə}R	&	*qat\tbr{ə}p	&	*waR\tbr{ə}j	&	*dal\tbr{ə}m	&	*\tbr{ə}n\tbr{ə}m\\	\midrule
\ili{Bobot}	&	{\hr}ma{\tl}met\tbr{a}n	&	{\hr}nih\tbr{a}n	&	{\hr}av-lok\tbr{a}	&	{\hr}yat\tbr{a}	&	{\hr}waw\tbr{a}t	&	{\hr}lal\tbr{a}	&	{\hr}num\\
\ili{Masiwang}	&	{\hr}meta{\tl}met\tbr{a}ŋ	&	{\hr}-nif\tbr{a}\cg{(n)}	&	{\hr}lok\tbr{a}	&	\cg{(ilan)}	&	{\hr}was\tbr{a}t	&	{\hr}lal\tbr{a}n	&	{\hr}noon\\	\midrule
\ili{Bati}	&	‹miat\tbr{a}n›	&	\cg{(‹nisikonina›)}	&		&		&		&	‹won\tbr{e}n›\\
\ili{Geser}	&	{\hr}met\tbr{a}n	&	\cg{(ŋisik)}	&	{\hr}roŋ\tbr{a}n	&	\cg{(barem)}	&	\cg{(tali)}	&	\cg{(lomin)}	&	{\hr}on\tbr{a}n\\
\ili{Gorom}	&	{\hr}met\tbr{a}n	&	\cg{(ŋisi)}	&	{\hr}roŋ\tbr{a}n	&	\cg{(barem)}	&	\cg{(tali)}	&	\cg{(lomin)}	&	{\hr}on\tbr{a}n\\
\ili{Watubela}	&	{\hr}maket\tbr{a}n	&	{\hr}nif\tbr{o}	&	{\hr}doŋ\tbr{a}n	&	{\hr}kat\tbr{a}f	&	\cg{(tali)}	&	\cg{(lomi)}	&	{\hr}on\tbr{o}n\\
\ili{Kesui}	&	‹met\tbr{e}n›	&	‹nif\tbr{oa}›	&		&		&		&		&	{\hr}on\tbr{a}m\\
\ili{Kowiai}	&	{\hr}maet\tbr{a}n	&	{\hr}rif\tbr{u}t	&	{\hr}fanoŋg\tbr{a}r	&	\cg{(barem)}	&	{\hr}war\tbr{a}s	&	\cg{(mora)}	&	\\	\midrule
\ili{Teor}	&	‹mit\tbr{e}n›	&	‹nif\tbr{i}n›	&		&		&		&		&	‹nem›\\
\ili{Kur}	&	‹mid\tbr{a}n›	&	{\hr}nip\tbr{a}n	&	{\hr}-nan\tbr{a}r	&		&		&		&	‹nèèn›\\	\midrule
\ili{Yamdena}	&	{\hr}me-met\tbr{a}m	&	{\hr}nif\tbr{a}-ny	&	{\hr}-deŋ\tbr{a}r	&	{\hr}daf\tbr{a}t\su{‡}	&	{\hr}war\tbr{a}s	&	{\hr}dal\tbr{a}m	&	{\hr}nem\\
\ili{Fordata}	&	{\hr}ŋarmer\tbr{a}n\su{†}	&	{\hr}nif\tbr{a}-n	&	{\hr}-ren\tbr{a}r	&	{\hr}raf\tbr{a}t\su{‡}	&	{\hr}war\tbr{a}t	&	{\hr}ral\tbr{a}-n	&	{\hr}ne\tbr{a}n\\
\ili{Kei}	&	{\hr}met-met\tbr{a}n	&	{\hr}nif\tbr{a}n	&	{\hr}den\tbr{a}r	&	{\hr}raat, raf\tbr{a}t\su{‡}	&	{\hr}war\tbr{a}t	&	{\hr}ra\tbr{a}-n	&	{\hr}ne\tbr{a}n\\
\ili{Onin}	&	\cg{(kamusa)}	&	{\hr}nif\tbr{a}n	&	\cg{(a-tanam)}	&	\cg{(idafi)}	&	{\hr}war\tbr{a}s	&	{\hr}nan\tbr{a}m	&	{\hr}nem\\
\ili{Sekar}	&	\cg{(kud{\tl}kuda)}	&	{\hr}nif\tbr{a}n	&	\cg{(tanam)}	&	\cg{(sirafakin)}	&	{\hr}war\tbr{a}s	&	{\hr}nan\tbr{a}m	&	{\hr}nem\\
Uruang.	&	\cg{(kom{\tl}koma)}	&	{\hr}nif\tbr{a}n	&	\cg{(a-tanam)}	&	\cg{(sanoŋ)}	&	{\hr}war\tbr{a}s	&	{\hr}nan\tbr{a}m	&	{\hr}nem\\
\ili{Arguni}	&	\cg{(ud{\tl}ude)}	&	\cg{(ifiye tutu)}	&	\cg{(a-tanem)}	&	\cg{(erár)}	&	{\hr}war\tbr{i}r	&	\cg{(ramon)}	&	{\hr}an\tbr{e}m\\
\ili{Bedoanas}	&	\cg{(amunda)}	&	\cg{(feʔe tutu)}	&		&		&	{\hr}war\tbr{i}r	&		&	{\hr}un\tbr{a}m\\
Erokw.	&	\cg{(amunda)}	&	\cg{(feʔe tutu)}	&	\cg{(tanam)}	&	\cg{(erár)}	&	{\hr}war\tbr{i}r	&	\cg{(-ramu)}	&	{\hr}any\tbr{a}m\\
\ili{Irarutu}	&	{\hr}grmutn	&	{\hr}rf\tbr{o}	&	{\hr}-f-nogr	&	{\hr}at\tbr{i}f-ro	&	{\hr}war\tbr{a}	&	\cg{(gan)}	&	\cg{(treso)}\\
\ili{Kuri}	&	\cg{(səsiːnə)}	&	{\hr}rɔb\tbr{o}	&	{\hr}itə-bə-ˈnɔgrə̆	&	{\hr}at\tbr{ɛ}βˈrɔ	&		&	\cg{(gurˈnɛ)}	&	\\
		\lspbottomrule
	\end{tabular}
		\begin{tablenotes}\tnsize
			\item[†]	{If \ili{Fordata} \it{ŋarmeran} ‘dark’ is from *ma-qitəm it has irregular *t > \it{r}.
								\ili{Fordata} has \it{ŋtoan} ‘black’.}
			\item[‡]	{If \ili{Kei} \it{rafat,} \ili{Fordata} \it{rafat}
								and \ili{Yamdena} \it{dafat} ‘(thatched) roof’ are from *qatəp, they
								have metathesis of the penultimate and final consonants.
								The source of the initial consonant is also unknown.}
		\end{tablenotes}
	\end{threeparttable}
\end{adjustbox}
\end{table}

Examples of *ə > **a in final syllables are given in \trf{11-04}
and examples of *a = **a in final syllables are given in \trf{11-05}.
Note in particular the reflexes of *ma-qitəm ‘black’ have *ə
> \it{a} in all languages.
(\ili{Teor} has subsequent **a > ‹e›,
‹a› /iC{\gap}{\#}; \srf{11-04}.) This sets STB apart from \tsc{\ili{Ambon}-Seram}
for which all known languages have *ə > \it{e} in this word,
including \tsc{Seti} \it{meten-a} ‘black’ (\srf{10-03-04}).
The reflexes in the \tsc{East Bomberai} languages (\ili{Arguni},
\ili{Sekar}, \ili{Onin}, \ili{Irarutu}, and \ili{Kuri}) are quite diverse,
including \it{a, e, i,} and {\0}.
Nonetheless, these diverse reflexes are found for both *ə
and *a, meaning we can posit a merger with subsequent split.
Similarly, \ili{Teor} and \ili{Kur} have *a/*ə > \it{i} for certain inalienably possessed words.
See \srf{11-04} for more details.
The reflexes of final *ə in *ənəm ‘six’ are also diverse
and do not conform to the usual patterns for final *ə.

\begin{table}\tcp{\tsc{Seram-Tanimbar-Bomberai} *a = ** a /σσ{\#}}{11-05}
\begin{adjustbox}{width=\textwidth}
	\begin{threeparttable}\st{2pt}
	\begin{tabular}{llllllll}\lsptoprule
local gl.	&	blood	&	road	&	house	&	moon	&	name	&	fish	&	root\\
PMP	&	*daR\tbr{a}q	&	*zal\tbr{a}n	&	*Rum\tbr{a}q	&	*bul\tbr{a}n	&	*ŋaj\tbr{a}n	&	*hik\tbr{a}n	&	*wak\tbr{a}[R/t]\\	\midrule
\ili{Bobot}	&	{\hr}law\tbr{a}	&	{\hr}lol\tbr{a}n	&	{\hr}um\tbr{a}	&	{\hr}vul\tbr{a}n	&	-nal\tbr{a}n	&	{\hr}i\tbr{a}n	&	{\hr}ai waʔ\tbr{a}t\\
\ili{Masiwang}	&	{\hr}las\tbr{a}	&	{\hr}lol\tbr{o}n	&	{\hr}sum\tbr{a}	&	{\hr}hul\tbr{a}n	&	{\hr}nal\tbr{a}-n	&	{\hr}i\tbr{a}n	&	\cg{(tamit)}\\	\midrule
\ili{Bati}	&	‹lal\tbr{a}i›	&	‹lā\tbr{ā}n›	&	‹lúm\tbr{e}›	&	‹wú\tbr{a}n›	&		&	‹ik\tbr{a}n›	&	‹ak\tbr{a}r›, aʔ\tbr{a}l\su{†}\\
\ili{Geser}-G.	&	{\hr}rar\tbr{a}	&	{\hr}lal\tbr{a}n	&	{\hr}rum\tbr{a}	&	{\hr}ul\tbr{a}n	&	{\hr}ŋas\tbr{a}n	&	{\hr}iʔ\tbr{a}n	&	{\hr}aʔ\tbr{a}r\\
\ili{Watubela}	&	{\hr}lal\tbr{a}k	&		&	{\hr}lum\tbr{a}k	&	{\hr}ul\tbr{a}n	&	{\hr}ŋah\tbr{a}n	&	{\hr}ik\tbr{a}n	&	{\hr}ak\tbr{a}l\\
\ili{Kesui}	&	‹lar\tbr{a}h›	&	{\hr}lar\tbr{a}n	&	{\hr}orum\tbr{a}	&	{\hr}wul\tbr{a}n	&		&	‹i′\tbr{a}n›	&	{\hr}ai ah\tbr{a}\\
\ili{Kowiai}	&	{\hr}rar\tbr{a}	&	{\hr}rar\tbr{a}n	&	\cg{(san)}	&	{\hr}ɸur\tbr{a}n	&	{\hr}nes\tbr{a}	&	\cg{(dón)}	&	{\hr}wa\tbr{a}r\\	\midrule
\ili{Teor}	&	‹lar\tbr{a}h›	&	\cg{(‹lagain›)}	&	\cg{(‹sarin›)}	&	‹phul\tbr{a}n›	&		&	‹ik\tbr{a}n›	&	‹wok\tbr{i}›\\
\ili{Kur}	&		&	\cg{(liŋain)}	&		&	{\hr}fu\tbr{a}n	&	{\hr}ŋa\tbr{a}-m	&	{\hr}ik\tbr{a}n	&	{\hr}wak\tbr{i}r\\	\midrule
\ili{Yamdena}	&	{\hr}dar\tbr{a}-ny	&	{\hr}dalnen-e\su{‡}	&	\cg{(das)}	&	{\hr}bul\tbr{a}n	&	{\hr}ŋar\tbr{a}-ny	&	{\hr}i\tbr{a}n	&	{\hr}wak\tbr{a}r-y\\
\ili{Fordata}	&	{\hr}lar\tbr{a}-n	&	\cg{(liŋaʔan)}	&	\cg{(rahan)}	&	{\hr}vul\tbr{a}n	&	{\hr}nar\tbr{a}-n	&	{\hr}i\tbr{a}n	&	{\hr}waʔ\tbr{a}r\\
\ili{Kei}	&	{\hr}lar	&	\cg{(laŋay)}	&	\cg{(raha)}	&	{\hr}vu\tbr{a}n	&	\cg{(mema-n)}	&	\cg{(vuʔut)}	&	{\hr}waʔ\tbr{a}r\\
\ili{Onin}	&	{\hr}rar\tbr{a}	&	{\hr}nam\tbr{a}n	&	{\hr}rum\tbr{a}	&	{\hr}pun\tbr{a}n	&	{\hr}gar\tbr{a}-r	&	\cg{(sai)}	&	{\hr}wak\tbr{e}r\\
\ili{Sekar}	&	{\hr}rar\tbr{a}	&	{\hr}nan\tbr{a}m	&	{\hr}rum\tbr{a}	&	{\hr}bun\tbr{a}n	&	{\hr}gar\tbr{a}-r	&	\cg{(sai)}	&	\cg{(sorsor)}\\
Uruang.	&	{\hr}nar\tbr{a}	&	{\hr}nenmatan	&	\cg{(kapala)}	&	{\hr}pun\tbr{a}n	&	{\hr}gar\tbr{a}n	&	\cg{(seir)}	&	{\hr}wak\tbr{a}r\\
\ili{Arguni}	&	{\hr}rar\tbr{e}	&	{\hr}rár\tbr{i}n	&	{\hr}rum\tbr{e}	&	{\hr}púr\tbr{i}n	&	{\hr}gád\tbr{i}n	&	\cg{(sair)}	&	\cg{(dareg)}\\
\ili{Bedoanas}	&	\cg{(rita)}	&	{\hr}rár\tbr{i}n	&	\cg{(windi)}	&	{\hr}úr\tbr{i}n	&	{\hr}gad\tbr{i}ŋ	&	\cg{(sai)}	&	\cg{(dara)}\\
\ili{Erokwanas}	&	\cg{(rita)}	&	{\hr}rór\tbr{i}n	&	\cg{(widi)}	&	{\hr}púr\tbr{i}n	&	{\hr}gad\tbr{i}n	&	\cg{(sai)}	&	\cg{(adara)}\\
\ili{Irarutu}	&	\cg{(wams)}	&	{\hr}rarn	&	\cg{(san)}	&	\cg{(seba)}	&	\cg{(no)}	&	\cg{(sum)}	&	{\hr}war\\
\ili{Kuri}	&	\cg{(wams)}	&	{\hr}rarn	&	\cg{(wɛˈnɔm)}	&	\cg{(sarabɛ)}	&	\cg{(inɔi)}	&	\cg{(isɔːm)}	&	{\hr}war\\
		\lspbottomrule
	\end{tabular}
		\begin{tablenotes}\tnsize
			\item[†]	{\ili{Bati} ‹akar› ‘root’ is from \citet[621]{Wallace1869}
								and \it{aʔal} is from \citet[127]{LoskiLoski1989}.}
		\end{tablenotes}
	\end{threeparttable}
\end{adjustbox}
\end{table}

Examples of *ə in final syllables in
addition to those in \trf{11-04} include: 

\begin{itemize}
	\item	*pus\tbr{ə}j ‘navel’ > \ili{Geser} \it{fus\tbr{a}l}; \ili{Gorom} \it{us\tbr{a}l}; \ili{Kowiai}
				\it{ɸus\tbr{a}r}; \ili{Irarutu} \it{fwir}; \ili{Kuri} \it{wuer-e},
				\ili{Uruangnirin} \it{fus\tbr{a}l}; \ili{Kei} \it{fuh\tbr{a}r}; \ili{Fordata} \it{fu\tbr{a}r};
				and \ili{Yamdena} \it{fus\tbr{a}r};
	\item	*qul\tbr{ə}j ‘worm,
				maggot’ > \ili{Fordata} \it{ul\tbr{a}r}
				and \ili{Bobot} \it{ur\tbr{a}t} ‘intestinal worm’;
	\item	*qabat\tbr{ə}d ‘sago grub’> \ili{Yamdena} \it{bat\tbr{a}r};
	\item	*tiR\tbr{ə}m ‘oyster’ > \ili{Bobot} \it{til\tbr{a}n} (irregular *R > \it{l});
	\item	*kəd\tbr{ə}ŋ ‘stand’ > \ili{Bobot},
				\ili{Masiwang} \it{kol\tbr{a}n}.
\end{itemize}

\begin{table}\tcp{\tsc{Seram-Tanimbar-Bomberai} *-ay}{11-06}
\begin{adjustbox}{width=\textwidth}
	\st{2pt}\begin{tabular}{lllllll}\lsptoprule
local gl.	&	sand	&	rice	&	chin, jaw	&	man	&	liver	&	die\\
PMP	&	*qən\tbr{ay}	&	*paj\tbr{ay}	&	*qaz\tbr{ay}	&	*maRuqan\tbr{ay}	&	*qat\tbr{ay}	&	*mat\tbr{ay}\\	\midrule
\ili{Bobot}	&	{\hr}yen\tbr{a}	&		&	{\hr}yal\tbr{a}-n	&	{\hr}muan\tbr{a}	&	{\hr}yat\tbr{a}-n	&	{\hr}mat\tbr{a}\\
\ili{Masiwang}	&	\cg{(bes)}	&	{\hr}fas\tbr{a}	&	{\hr}yal\tbr{a}-n	&	{\hr}mesan\tbr{e}m, masan\tbr{a}n	&	{\hr}yat\tbr{a}-n	&	{\hr}mat\tbr{a}\\	\midrule
\ili{Bati}	&		&	‹faas\tbr{í}›	&		&	‹bulan\tbr{a}› ‘husband’	&		&	\\
\ili{Geser}	&	{\hr}en\tbr{a}	&	{\hr}fas\tbr{a}	&	{\hr}ar	&	{\hr}uran\tbr{a}	&	{\hr}at	&	{\hr}mat\tbr{a}\\
\ili{Gorom}	&	{\hr}en\tbr{a}	&	{\hr}has\tbr{a}	&	{\hr}ar	&	{\hr}uran\tbr{a}	&	{\hr}at	&	{\hr}mat\tbr{a}\\
\ili{Watubela}	&	{\hr}ken\tbr{a} wo	&		&	\cg{(kakelan)}	&	{\hr}maŋkan\tbr{a}	&	\cg{(lamno-)}	&	{\hr}-mat\tbr{a}\\
\ili{Kesui}	&		&	{\hr}fah\tbr{a}	&		&		&		&	\\
\ili{Kowiai}	&	{\hr}en\tbr{a}	&	{\hr}ɸas\tbr{a}	&	{\hr}lar\tbr{a}ʔai	&	{\hr}muan\tbr{a}, muruan\tbr{a}	&	{\hr}lat\tbr{a}	&	{\hr}na-mat\tbr{a}\\	\midrule
\ili{Teor}	&		&		&		&	‹meránn\tbr{a}›	&		&	\\
\ili{Kur}	&		&	{\hr}pas\tbr{a}	&		&	{\hr}manran\tbr{a}	&		&	{\hr}-mat\tbr{a}/-mat\\	\midrule
\ili{Yamdena}	&	\cg{(knryai)}	&	{\hr}fas-e	&	\cg{(keka-ny)}	&	{\hr}merwan-e	&	{\hr}at\tbr{a}-ŋw	&	{\hr}n-mat\\
\ili{Fordata}	&	\cg{(ŋuur)}	&	\cg{(wanat)}	&	\cg{(demi-n)}	&	{\hr}bran\tbr{a}	&	{\hr}yat\tbr{a}-n	&	{\hr}n-mat\tbr{a}\\
\ili{Kei}	&	\cg{(ŋuur)}	&	\cg{(kokat)}	&	\cg{(kakean)}	&	{\hr}bran	&	{\hr}yat\tbr{a}-n	&	{\hr}n-mat\\
\ili{Onin}	&	\cg{(igai)}	&	{\hr}pas\tbr{a}	&	\cg{(keken)}	&	{\hr}morar\tbr{a}-r	&	{\hr}yat\tbr{a}-n	&	{\hr}a-mat\tbr{a}\\
\ili{Sekar}	&	\cg{(kidai)}	&	{\hr}pas\tbr{a}	&	\cg{(kaken)}	&	{\hr}morar\tbr{a}	&	\cg{(nawan)}	&	\cg{(warag)}\\
Uruang.	&	\cg{(idei)}	&	{\hr}fas\tbr{a}	&	\cg{(keken)}	&	{\hr}marwan\tbr{a}	&	{\hr}lat\tbr{a}-n ‘lungs’	&	{\hr}mata puan\\
\ili{Arguni}	&	\cg{(ain)}	&	{\hr}pwas\tbr{e}n	&	\cg{(aure-m)}	&	{\hr}maran	&		&	{\hr}ga-mat\\
\ili{Bedoanas}	&	\cg{(ain)}	&	{\hr}pas\tbr{a}	&		&	{\hr}maran, manyan\tbr{a}	&	\cg{(tapoti)}	&	{\hr}mu-mat\\
\ili{Erokwanas}	&	\cg{(ain)}	&	{\hr}pas\tbr{a}	&		&	{\hr}maran	&	\cg{(wade)}	&	{\hr}i-m\<y\>at\\
\ili{Irarutu}	&	{\hr}eny\tbr{e}fu	&	{\hr}fas	&	\cg{(kke)}	&	{\hr}mran	&	{\hr}t\tbr{i}	&	{\hr}n-mat\\
	&	[eny\tbr{e}fu]	&	[fas\tbr{ə}]	&	[kəke]	&	[məran\tbr{ə}]	&	[t\tbr{i}]	&	[nəmat\tbr{ə}]\\
\ili{Kuri}	&	{\hr}en\tbr{ə}fafo	&		&	\cg{(‹ˈae›)}	&	{\hr}bran	&		&	{\hr}inə-mat\\
		\lspbottomrule
	\end{tabular}
\end{adjustbox}
\end{table}

Examples of final vowel-glide sequences are given in \trf{11-06}
for *-ay and \trf{11-07} for *-aw.
(\ili{Arguni}, \ili{Bedoanas}, and \ili{Erokwanas} are not given in \trf{11-07}
as no reflexes of the PMP etyma in this table have been identified.)
An additional example of final *-ay does not consistently have *-ay > \it{a}.
This is *bəRsay ‘oar’ > \ili{Arguni} \it{pwores},
\ili{Bedoanas} \it{oras}, \ili{Erokwanas} \it{poras},
\ili{Irarutu} \it{eʤe}, \ili{Kuri} \it{boye},
\ili{Onin}, \ili{Uruangnirin} \it{pesa}, \ili{Sekar} \it{besa},
\ili{Kei} \it{vehe}, \ili{Fordata} \it{vahi},
\ili{Yamdena} \it{bese}.

\begin{table}[t]\centering\tcp{\tsc{Seram-Tanimbar-Bomberai} *-aw}{11-07}
\begin{adjustbox}{width=\textwidth}
	\begin{threeparttable}\st{2pt}
	\begin{tabular}{lllllll}\lsptoprule
local gl.	&	day, sun	&	walk, go	&	mouse, rat	&	steal	&	rafter	&	sister (m.s.)\\
PMP	&	*qaləj\tbr{aw}	&	*(ma)-pan\tbr{aw}	&	*lab\tbr{aw}	&	*nak\tbr{aw}	&	*kas\tbr{aw}	&	*bət\tbr{aw}\\	\midrule
\ili{Bobot}	&	{\hr}li\tbr{a} matan	&	\cg{(ko)}	&	{\hr}milav\tbr{a}	&	{\hr}ma-na\tbr{a}	&	{\hr}as\tbr{a}-	&	{\hr}bot\tbr{a}-n\\
\ili{Masiwang}	&	\cg{(roman)}	&	\cg{(uha)}	&	{\hr}ilah\tbr{a}	&	\cg{(kah leit)}	&	{\hr}as\tbr{a}	&	\cg{(lasa)}\\	\midrule
\ili{Bati}	&	‹wol\tbr{eh}›	&	\cg{(‹ketango›)}	&	\cg{(‹kaúfei›)}	&		&		&	\\
\ili{Geser}	&	{\hr}ol\tbr{a}	&	\cg{(taŋgi)}	&	\cg{(kalufa)}	&	\cg{(fakaleus)}	&		&	\\
\ili{Gorom}	&	{\hr}ol\tbr{a}	&	\cg{(tagi)}	&	\cg{(aruha)}	&	\cg{(haleus)}	&	{\hr}as\tbr{a}r	&	\\
\ili{Watubela}	&		&	{\hr}fan\tbr{a}	&		&	{\hr}komanak\tbr{a}	&		&	\\
\ili{Kesui}	&	‹ol\tbr{ēr}›	&	‹fan\tbr{ów}›	&	\cg{(arofa)}	&		&		&	\\
\ili{Kowiai}	&	{\hr}or\tbr{a}	&	{\hr}seɸan\tbr{a}	&	\cg{(asaɸai)}	&	\cg{(na-ɸareʔeus)}	&	{\hr}ʔas\tbr{a}	&	{\hr}ɸut\tbr{a} ‘ySi’\\	\midrule
\ili{Teor}	&	‹lew›	&	\cg{(‹takek›)}	&	\cg{(‹fudarúa›)}	&		&		&	\\
\ili{Kur}	&	{\hr}loŋ, ‹lè›	&	\cg{(i-kek)}	&		&		&		&	\\	\midrule
\ili{Yamdena}	&	{\hr}ler-e	&	{\hr}na-mpan	&	\cg{(kawat)}	&	\cg{(na-mnaŋ)}	&	{\hr}as-e	&	{\hr}abut\\
\ili{Fordata}	&	{\hr}ler\tbr{a}	&	{\hr}n-ban\tbr{a}	&	\cg{(manovan)}	&	\cg{(n-bori)}	&	{\hr}ah\tbr{a}	&	\cg{(ura-n)}\\
\ili{Kei}	&	{\hr}ler	&	\cg{(ba)}\su{‡}	&	\cg{(keru)}	&	\cg{(n-bor)}	&	{\hr}ah\tbr{a}	&	\cg{(ura-n)}\\
\ili{Onin}	&	{\hr}rer\tbr{a}	&	{\hr}a-ban\tbr{a}	&	\cg{(yawar)}	&	\cg{(nənena)}	&		&	\\
\ili{Sekar}	&	{\hr}rer\tbr{a}	&	{\hr}a-pan\tbr{a}, fan\tbr{a}	&	\cg{(kufiata)}	&	\cg{(afnage)}	&		&	\\
Uruang. &	{\hr}ner\tbr{a}	&	{\hr}a-ban\tbr{a}	&	\cg{(yawat)}	&	\cg{(amnaŋgawai)}	&	\cg{(lepir)}	&	\\
\ili{Irarutu}\su{†}	&	{\hr}re, rre	&	{\hr}bane	&	\cg{(sfe)}	&	\cg{(nəbəna)}	&	\cg{(sasuirə)}	&	\\
\ili{Kuri}	&	{\hr}rəˈrɔːr	&	\cg{(itə-βa)}\su{‡}	&	\cg{(kɔkbɛru)}	&		&		&	\\
		\lspbottomrule
	\end{tabular}
		\begin{tablenotes}\tnsize
			\item[†]	{\ili{Irarutu} \it{sasuirə} ‘rafter’,
								\it{bane} ‘walk’, and \it{nəbana} ‘steal’ are from \citet{Voorhoeve1995b}.
								The final vowel in \it{bane} ‘walk’ is probably epenthetic.}
			\item[‡]	{\ili{Kei} \it{ba}
								and \ili{Kuri} \it{itə-βa} are probably from PMP *ba ‘at,
								on, in, to’ or PCEMP **ba ‘go’.}
		\end{tablenotes}
	\end{threeparttable}
\end{adjustbox}
\end{table}

A feature of \ili{Yamdena} grammar that seems to have clouded some
of the comparative literature
is that the citation form in word lists often ends in \it{e{\#}},
masking underlying \it{a} which surfaces in other forms.
This final \it{-e} is analysed by \citet{MettlerMettler1997}
as a determiner suffix.
Examples include *qalima > \it{lim-e} ‘hand/arm’,
\it{lima-ŋw} ‘my arm’, and *qatay > \it{at-e} ‘liver’,
\it{ata-ŋw} ‘my liver’.\footnote{
		 Formally similar nominal markers
		also occur in the languages of \tsc{Southwest Maluku} (\srf{04-06}) and \tsc{\ili{Ambon}-Seram}.
		\citet{Stresemann1927} posited a support vowel \it{-e} (\srf{07-01-01}) for Sub-\ili{Ambon}.
		%For \ili{Leti}, \citet[159f]{vanEngelenhoven2004} describes a definiteness
		%marker \it{-e} which only occurs on nouns with final /a/,
		%in which case it replaces this final vowel.
		}
\ili{Yamdena} usually
has loss of *ə/*a/*-ay/-*aw in final open syllables.
